\chapter{Results and Post-Processing}
\label{ch:results}

This chapter presents the integral fluid-dynamic results of the final high-pressure, low-temperature simulations. As the target application is an excess-flow valve (EFV) for hydrogen-vehicle refuelling, \emph{hydrogen is the primary working fluid}. Nitrogen is investigated as a supplementary industrial gas and its numerical response is discussed after the hydrogen findings. The full-factorial study comprises four obturator geometries, three inlet velocities ($5$, $10$ and $15\,\mathrm{m\,s^{-1}}$), and the two working gases.

The final simulation archive is identified with the nominal reference state $437\,\mathrm{bar}$ and $233\,\mathrm{K}$ (approximately $-40\,^{\circ}\mathrm{C}$). These values identify the chosen operating condition; the constant-density solver does not itself solve a real-gas equation of state. All values in the tables below are computed OpenFOAM responses, not experimental measurements.

\section{Structure of OpenFOAM Post-Processing Data}
\label{sec:postprocessing_structure}

Each calculation writes five sets of numerical outputs: area-averaged kinematic pressure at the inlet and outlet, integrated inlet velocity, pressure and viscous loads on the obturator, and wall-distance data. The corresponding function-object locations are \texttt{inlet\_pressure\_avg}, \texttt{outlet\_pressure\_avg}, \texttt{inlet\_volume\_flux}, \texttt{forceCoeffs\_Obt} and \texttt{yPlus1}. Although the force output folder includes the term force coefficients, the active OpenFOAM function-object type is \texttt{forces}, not \texttt{forceCoeffs}.

The solved pressure field is \emph{kinematic} pressure, expressed in $\mathrm{m^2\,s^{-2}}$. The dimensional pressure loss and axial force component are evaluated using
\begin{equation}
\Delta p=\rho_{\mathrm{ref}}\left(P_{\mathrm{in}}-P_{\mathrm{out}}\right),
\qquad F_x=F_{p,x}+F_{\nu,x},
\label{eq:processed_responses}
\end{equation}
where $P_{\mathrm{in}}$ and $P_{\mathrm{out}}$ are area averages and $F_{p,x}$ and $F_{\nu,x}$ are the pressure and viscous contributions to the surface-integrated force. The output \texttt{forces.dat} already uses the gas-specific \texttt{rhoInf} and is expressed in newtons; it must \emph{not} be multiplied by density a second time. The global positive-$x$ force is distinct from a mechanically verified closing-direction force.

The submitted dictionaries give $\rho_{\mathrm{H}_2}=33.77\,\mathrm{kg\,m^{-3}}$, $\nu_{\mathrm{H}_2}=2.709\times10^{-7}\,\mathrm{m^2\,s^{-1}}$, $\rho_{\mathrm{N}_2}=481.352\,\mathrm{kg\,m^{-3}}$ and $\nu_{\mathrm{N}_2}=7.932\times10^{-8}\,\mathrm{m^2\,s^{-1}}$. The first quantity for each gas is the active force-reference density and the second is the viscosity used by the incompressible solver. The thermodynamic interpretation of these reference settings is discussed in Section~\ref{sec:doe_highpressure_state}.

\section{Data Reduction and Terminal Numerical Behaviour}
\label{sec:plotly_dashboard}

The \texttt{simpleFoam} output index represents solver iteration, not elapsed physical time. All 24 supplied configurations have pressure and force histories reaching iteration 3000. To apply a reproducible comparison to each case, the final 301 samples (indices 2700--3000 inclusive) were extracted. For a numerical response $R_i$, the arithmetic mean and the coefficient of variation are
\begin{equation}
\overline{R}=\frac{1}{N}\sum_{i=1}^{N}R_i,
\qquad
\mathrm{CV}(R)=
\frac{\sqrt{\frac{1}{N}\sum_{i=1}^{N}(R_i-\overline{R})^2}}
{|\overline{R}|}\times100\%,\qquad N=301.
\label{eq:terminal_average}
\end{equation}
This statistic measures variation within the selected terminal interval. Reaching the last iteration or observing a small CV does not independently prove residual convergence, mesh independence, or experimental validation \cite{moukalled2016,openfoam_user}.

For hydrogen, the terminal-window CV ranges from approximately $0.83$--$10.25\%$ for pressure loss and $0.65$--$12.46\%$ for axial force. The largest variations occur for obturator B at $15\,\mathrm{m\,s^{-1}}$. In that case, comparing the final and initial 100 samples within the terminal window gives changes of approximately $11.6\%$ in mean pressure loss and $15.1\%$ in mean force. Obturator B at $10\,\mathrm{m\,s^{-1}}$ also exhibits appreciable variation. These cases are therefore reported as \emph{terminal-window estimates} rather than proven time-independent limits.

\subsection{Inlet flow consistency and mass-flow interpretation}

The integrated inlet area is approximately $A_{\mathrm{in,mesh}}=4.658925\times10^{-5}\,\mathrm{m^2}$. The inlet volumetric flow rates are approximately $2.32946\times10^{-4}$, $4.65893\times10^{-4}$ and $6.98839\times10^{-4}\,\mathrm{m^3\,s^{-1}}$ at inlet velocities of 5, 10 and $15\,\mathrm{m\,s^{-1}}$, respectively. These are based on the discretized patch rather than the ideal CAD inlet area. For the constant-density model, the reference mass-flow rate is
\begin{equation}
\dot m_{\mathrm{ref}}=\rho_{\mathrm{ref}} Q
=\rho_{\mathrm{ref}}U_{\mathrm{in}}A_{\mathrm{in,mesh}}.
\label{eq:mass_flow_highp}
\end{equation}
The hydrogen mass-flow rates at the three levels are therefore approximately $0.00787$, $0.01573$ and $0.02360\,\mathrm{kg\,s^{-1}}$. At the same imposed velocities, the corresponding nitrogen mass-flow rates are approximately $0.1121$, $0.2243$ and $0.3364\,\mathrm{kg\,s^{-1}}$. Thus equal-velocity gas comparisons are not equal-mass-flow refuelling comparisons.

\section{Hydrogen: Primary High-Pressure CFD Results}
\label{sec:fluid_trends}
\label{sec:hydrogen_highp_results}

Tables~\ref{tab:hydrogen_pressure} and~\ref{tab:hydrogen_force} give the hydrogen pressure losses and total positive-$x$ obturator forces obtained from the common terminal interval. Pressure differences are reported in kilopascals, while the force outputs are already dimensional newtons.

\begin{table}[htbp]
\centering
\begin{tabular}{@{}rcccc@{}}
\toprule
$U_{\mathrm{in}}$ [$\mathrm{m\,s^{-1}}$] & \textbf{Obt.\ A} & \textbf{Obt.\ B} & \textbf{Obt.\ C} & \textbf{Obt.\ D}\\
\midrule
5  & 33.36 & 3.45 & 9.27 & 1.72\\
10 & 132.57 & 15.68 & 37.60 & 6.86\\
15 & 297.61 & 34.81 & 83.73 & 15.55\\
\bottomrule
\end{tabular}
\caption{Hydrogen: mean inlet-to-outlet pressure loss $\overline{\Delta p}$ [kPa] at the high-pressure reference state, solver iterations 2700--3000.}
\label{tab:hydrogen_pressure}
\end{table}

\begin{table}[htbp]
\centering
\begin{tabular}{@{}rcccc@{}}
\toprule
$U_{\mathrm{in}}$ [$\mathrm{m\,s^{-1}}$] & \textbf{Obt.\ A} & \textbf{Obt.\ B} & \textbf{Obt.\ C} & \textbf{Obt.\ D}\\
\midrule
5  & 5.616 & 0.457 & 1.468 & 0.165\\
10 & 22.344 & 2.094 & 5.947 & 0.674\\
15 & 50.089 & 4.768 & 13.281 & 1.531\\
\bottomrule
\end{tabular}
\caption{Hydrogen: mean axial hydrodynamic force component $\overline{F}_x$ [N] at the high-pressure reference state, solver iterations 2700--3000. Positive force follows the computational $+x$ axis.}
\label{tab:hydrogen_force}
\end{table}

At each investigated inlet velocity, both response magnitudes follow the order A, C, B, D from largest to smallest. This trend is qualitatively consistent with the nominal total hole areas: a more restrictive passage produces stronger acceleration, pressure redistribution, and integrated loading. The precise force nonetheless depends on the three-dimensional pressure and shear distributions and not solely on the nominal area \cite{white2016,crane410}.

At $15\,\mathrm{m\,s^{-1}}$, obturator A gives approximately $297.6\,\mathrm{kPa}$ ($2.98\,\mathrm{bar}$) pressure loss and $50.09\,\mathrm{N}$ axial force; obturator D gives approximately $15.55\,\mathrm{kPa}$ ($0.156\,\mathrm{bar}$) and $1.53\,\mathrm{N}$. Obturator A therefore imposes approximately 19 times the pressure loss and generates approximately 33 times the axial load of obturator D at this condition. The trade-off between acceptable normal-flow loss and hydrodynamic load relevant to actuation is a central consideration for the company-facing EFV design.

Pressure loss and force increase strongly with inlet speed for every configuration. Their velocity dependence is broadly consistent with the importance of dynamic pressure, proportional to $\rho U_{\mathrm{in}}^2$, but the terminal fluctuations in some cases make a precise fitted exponent inappropriate without further numerical analysis.

\section{Nitrogen: Secondary Industrial-Gas Comparison}
\label{sec:nitrogen_comparison}

Hydrogen is the actual refuelling fuel and retains priority in the design interpretation. Nitrogen was included as a second working fluid because it is a familiar industrial gas in hydrogen-system purging and inerting procedures; its inclusion broadens the physical-property comparison without making it a substitute for hydrogen in fuel delivery \cite{h2tools_purging,doe_hydrogen_practices}. Nitrogen and hydrogen differ substantially in density, viscosity and resulting Reynolds number at the same imposed velocity.

The nitrogen results are reported in Tables~\ref{tab:nitrogen_pressure} and~\ref{tab:nitrogen_force}. One data-integrity qualification is necessary. In the supplied archive, the nitrogen case with obturator A at $5\,\mathrm{m\,s^{-1}}$ has force, pressure, inlet-flow and $y^+$ histories \emph{identical at the byte level} to the corresponding hydrogen outputs, despite using different physical-property settings in the case dictionaries. Therefore, its numerical response cannot be considered independently established; the tables identify that entry as not independently assessable.

\begin{table}[htbp]
\centering
\begin{tabular}{@{}rcccc@{}}
\toprule
$U_{\mathrm{in}}$ [$\mathrm{m\,s^{-1}}$] & \textbf{Obt.\ A} & \textbf{Obt.\ B} & \textbf{Obt.\ C} & \textbf{Obt.\ D}\\
\midrule
5 & n.a. & 52.61 & 130.74 & 24.90\\
10 & 1877.86 & 206.79 & 526.50 & 99.30\\
15 & 4225.06 & 466.30 & 1178.52 & 224.35\\
\bottomrule
\end{tabular}
\caption{Nitrogen: mean pressure loss $\overline{\Delta p}$ [kPa] at the high-pressure reference state, iterations 2700--3000. The 5 m/s obturator-A output is excluded due to duplicated source histories.}
\label{tab:nitrogen_pressure}
\end{table}

\begin{table}[htbp]
\centering
\begin{tabular}{@{}rcccc@{}}
\toprule
$U_{\mathrm{in}}$ [$\mathrm{m\,s^{-1}}$] & \textbf{Obt.\ A} & \textbf{Obt.\ B} & \textbf{Obt.\ C} & \textbf{Obt.\ D}\\
\midrule
5 & n.a. & 7.125 & 20.732 & 2.383\\
10 & 316.116 & 28.051 & 83.294 & 9.844\\
15 & 711.635 & 62.492 & 187.091 & 21.799\\
\bottomrule
\end{tabular}
\caption{Nitrogen: mean total axial force $\overline{F}_x$ [N] at the high-pressure reference state, iterations 2700--3000. The 5 m/s obturator-A output is excluded due to duplicated source histories.}
\label{tab:nitrogen_force}
\end{table}

The magnitude of the nitrogen loads is generally much greater than that of hydrogen at equal inlet velocity and geometry. At $15\,\mathrm{m\,s^{-1}}$ with obturator A, nitrogen gives approximately $4.23\,\mathrm{MPa}$ and $711.6\,\mathrm{N}$, compared with $0.298\,\mathrm{MPa}$ and $50.1\,\mathrm{N}$ for hydrogen. This difference is largely associated with the much higher nitrogen reference density, together with changes in the Reynolds number. It must not be interpreted as a comparison at equal mass-flow demand or as direct evidence about hydrogen refuelling performance.

\section{Near-Wall Diagnostics and Numerical Interpretation}
\label{sec:near_wall_verification}

The final post-processing outputs include minimum, maximum and average wall-distance values for the selected patches. At iteration 3000, the maximum $y^+$ reported on the obturator ranges from approximately 156 to 1256 across the hydrogen cases. The nitrogen cases exhibit larger maxima, exceeding 4000 in the most restrictive and highest-velocity combination. These values warrant careful interpretation of near-wall turbulence modelling, viscous loads and wall-layer resolution: they cannot be used as evidence that all surfaces are resolved with $y^+\approx1$ \cite{menter1994,openfoam_user}.

The results are numerical terminal-window estimates, not experimental validation of EFV trip behavior. Residual histories, outlet-flux balance, mesh independence and local passage velocity are distinct verification quantities. For hydrogen, the greatest terminal fluctuation occurs in obturator B at the higher velocities. For nitrogen, the maximum pressure loss of approximately $4.23\,\mathrm{MPa}$ represents about $9.7\%$ of the nominal $43.7\,\mathrm{MPa}$ reference pressure, making constant density more consequential for the secondary gas comparison. Even when an overall pressure ratio is moderate, local Mach number may increase in the narrow passages \cite{anderson2017}.

\section{Comparison Between Reference-State and Dispensing Conditions}
\label{sec:results_operating_states}

The earlier ambient-condition simulations provide a useful baseline, but the final high-pressure and low-temperature campaign is more representative of the specified EFV location downstream of the pre-cooler. The two states must be treated separately. In particular, changing the reference density in an already computed force history does not account for a different kinematic viscosity, Reynolds number or flow structure, and equal inlet velocity does not imply equal mass flow when density changes.

The application operating point must also be distinguished from the transient refuelling history governed by station-level controls. SAE~J2601 establishes light-duty gaseous-hydrogen filling protocols and delivery-temperature categories; the CFD cases instead describe flow past fixed obturators at specific reference properties \cite{sae_j2601}. The resulting axial fluid force is a component of, rather than a complete solution to, the valve's mechanical actuation problem.

\section{Implications for Excess-Flow Valve Design}
\label{sec:results_efv_implications}

The principal hydrogen result is a consistent trade-off: smaller effective passages increase both hydraulic resistance and positive-$x$ obturator loading. Therefore the configuration yielding the greatest force is not automatically the best valve, because the pressure-loss requirement during normal flow must also be met. An engineering choice requires the permissible loss, mechanical force balance, motion direction, spring properties and expected excess-flow conditions (Section~\ref{sec:efv_principle}).

The nitrogen study supports the comparison of fluid-property effects but does not change the hydrogen-first objective of the work. The tabulated CFD responses provide the numerical foundation for a design comparison tool, which is described in Chapter~\ref{ch:selection_tool}. Implementing and testing that software is a separate final development activity, not an additional round of flow simulations.
